\section{Motivation and scope}

Feature scaling sets relative learning rates as well as relative penalties.
For a rare binary feature, scaling by its standard deviation removes the
small marginal variance from the population Hessian, but it also increases
the magnitude of individual observations. Dependence with common features
can further change the learning trajectory. We ask when normalization helps
learn rare features under a finite training budget, and when stochastic
updates or dependence offset that benefit.

The connection between early stopping and regularization is established:
\citet{ali2019early} compare the risks of least-squares gradient flow and
ridge regression. Optimization geometry also determines implicit bias in
underdetermined models \citep{gunasekar2018geometry}, and preconditioning can
help or hurt generalization \citep{amari2020preconditioning}. Our starting
point is full-rank linear regression, where convergence to the unregularized
optimum gives ordinary least squares. The question concerns finite-time
learning and stochastic error; distinct interpolating solutions are a
possible later extension.

Rare predictive features already motivate AdaGrad
\citep{duchi2011adaptive}. More recently, \citet{kunstner2024imbalance}
connect heavy-tailed output-class imbalance to optimizer performance, while
\citet{sreckovic2025batch} show the importance of batch size and careful
tuning in language-model comparisons. Correlated input-feature rarity must
therefore contribute more than another example of slow learning under
imbalance. The candidate contribution is a joint account of dependence,
feature scaling, batch size, and prediction for rare subgroups. The exact
relationship to existing results remains to be established.

The calculations below extend the elementary two-feature flow and scalar
stability results to a fixed common block, finite-update error moments, and
batch-rarity stability under explicit population assumptions. Their relation
to the closest prior theorems remains a research objective. The explicit-penalty
analysis in the appendices supplies complementary baselines, extending the
normalization questions of \citet{larsson2025normalization}. Finite-sample
weighted-lasso distributions and polyhedral selection geometry already have
general treatments \citep{ewald2020distribution,lee2016exact}.
